\documentclass[11pt,a4paper]{article}
\usepackage{fontspec,amsmath,amssymb,amsthm,microtype}
\setmainfont{lmroman10-regular.otf}[BoldFont=lmroman10-bold.otf,ItalicFont=lmroman10-italic.otf,BoldItalicFont=lmroman10-bolditalic.otf]
\usepackage[margin=25mm]{geometry}
\usepackage[unicode,hidelinks]{hyperref}
\setlength{\emergencystretch}{3em}
\setlength{\parskip}{0.4em}
\setlength{\parindent}{0pt}
\allowdisplaybreaks[2]
\raggedbottom
\title{\texorpdfstring{Actual MYPGD bound about the original marginal center}{Actual MYPGD bound about the original marginal center}}
\author{Omer Deniz Akyildiz\\[0.4em]\normalsize Prepared with AI assistance}
\date{}
\begin{document}\maketitle
Draft mathematical exposition; these notes may contain mistakes.


The original-center MYPGD theorem adds a proved noise-removal radius to the corresponding MYIPLA estimate. All laws are identified with their actual Gaussian recursions. The resulting five contributions are the comparison radius, finite smoothing displacement, equilibrium spread, initial relaxation and smoothed Euler discretization.




Under the MYIPLA primitive assumptions, also assume \(m\le L_\eta=L_1+\eta^{-1}\). The algorithm uses genuine latent-only Gaussian updates; \(N>0\), \(m,\eta,\gamma>0\), and \(\gamma\le2/(m+L_\eta)\).


\section{\texorpdfstring{Identify actual algorithm laws in the comparison}{Identify actual algorithm laws in the comparison}}


Rewrite both output laws as the Gaussian Markov recursions of the shared smoothed drift. The initialization pushforward has finite second moment because its representative belongs to \(L^2\). The iterated noise-removal theorem now proves finite transport distance and \(W_2(\rho_n^{\rm PGD},\rho_n^{\rm IPLA})^2\le2d_\theta/(Nm)\).





Both chains have the same initial pushforward and deterministic smoothed forward map. The Gaussian path law and its independence identify them with the two Markov transitions analyzed earlier. The noise-removal theorem therefore proves a finite transport comparison for their actual output laws. Its squared upper bound is \(2d_\theta/(Nm)\), with every second moment established from the recursion rather than assumed from the desired rate.


\section{\texorpdfstring{Transfer parameter RMS through a nonexpansive projection}{Transfer parameter RMS through a nonexpansive projection}}


The parameter map has operator norm at most one. The projected RMS transport inequality yields \(\operatorname{RMS}_{\rm PGD}(\theta)\le W_2(\rho_n^{\rm PGD},\rho_n^{\rm IPLA})+\operatorname{RMS}_{\rm IPLA}(\theta)\). MYIPLA finite moments and the finite comparison imply the MYPGD moment is finite before passing to real square-root integrals. Taking the square root of the comparison gives the added radius \(\sqrt{2d_\theta/(Nm)}\).




Finite second moments of both algorithm laws give finite transport. Projecting a coupling onto the common parameter, whose operator norm is at most one, gives \[R_{\mathrm{MYPGD}}(\theta_*)\le R_{\mathrm{MYIPLA}}(\theta_*)+W_2(\rho_n^{\mathrm{MYPGD}},\rho_n^{\mathrm{MYIPLA}}).\] The actual-law comparison bounds the squared last term by \(2d_\theta/(Nm)\); nonnegative square roots therefore bound it by \(\sqrt{2d_\theta/(Nm)}\). The final theorem inserts the already proved original-center MYIPLA estimate, retaining its smoothing, equilibrium and discretization contributions.





Project a coupling of the two algorithm laws onto the common parameter. The projection cannot increase its quadratic cost. Adding the MYIPLA parameter error about the desired center and using Minkowski gives the MYPGD RMS bound. Taking the infimum over couplings gives the Wasserstein term, and its squared comparison estimate gives \(\sqrt{2d_\theta/(Nm)}\). Finite moments justify every square-root expectation involved.


\section{\texorpdfstring{Insert the proved original-\allowbreak{}center MYIPLA theorem}{Insert the proved original-center MYIPLA theorem}}


Apply the parameter-projection transport inequality to the two actual laws, then use the proved MYIPLA original-center theorem. With \(L_\eta=L_1+1/\eta\) and the assumptions above,


\[
\begin{aligned}
\left(\mathbb E\|\theta^{\mathrm{PGD}}_n-\theta_\star\|^2\right)^{1/2}
\le{}&\sqrt{\frac{2d_\theta}{Nm}}
 +\sqrt{\frac{\eta K^2}{m}}+\sqrt{\frac{d_\theta}{Nm}}\\
&+(1-m\gamma)^nW_2(\rho_0,\Pi_{N,\eta})\\
&+\frac{33}{20}\frac{L_\eta}{m}
 \sqrt{d_x+\frac{d_\theta}{N}}\sqrt\gamma.
\end{aligned}
\]


Every term is finite under the primitive assumptions, and every law in the statement is the actual defined algorithm law. This is a genuine finite-particle replacement for Appendix B’s Theorem B.1. It has a Gibbs initial-law transient and an explicit conservative comparison term; it does not reproduce the paper’s mean-field metric, constants, or stronger Taylor smoothing expression.





The MYIPLA bound already refers to the original marginal maximizer and includes its finite smoothing displacement. Adding the comparison radius leaves its other terms unchanged. The initial transient remains measured against the smoothed Gibbs reference, and the discretization coefficient retains \(L_1+1/\eta\). The result is therefore a finite-particle bound with explicit dependence on the smoothing scale, distinct from a stronger path-supremum or mean-field estimate.


\end{document}
